\section{Trusted Execution Environments}
\label{thesis:intro-cc}

With the term \acl{TEE} we refer to a set of \emph{hardware} extensions that can
be used to protect an application while executing within an (untrusted) system.
While there are technical differences between different architectures, all
\acp{TEE} aim to provide two main features:
%
\begin{enumerate}
    \item \emph{Hardware-based isolation} of logical resources to give
    confidentiality and integrity guarantees to application code and data. This
    is enforced in various ways depending on the specific \ac{TEE} design, but
    in general it involves ISA extensions to provide instructions for
    enabling/disabling protections and an access control mechanism to prevent
    illegal access to protected memory and CPU registers. The latter can be
    purely logical (e.g., allow or deny access based on the current program
    counter), cryptographic (e.g., encrypt and authenticate \ac{TEE} memory), or
    both. Some architectures may also bring additional elements to provide e.g.,
    some form of protection against side channels and physical attacks.
    \item \emph{Remote attestation}, a procedure to verify that an application
    is indeed running within a \ac{TEE} as expected. Here, a remote
    \emph{verifier}, which can be the application owner, can obtain a piece of
    \emph{evidence} containing information about the platform itself, such as
    firmware and microcode version, as well as information about the running
    workload, such as a measurement of the memory layout of the application at
    boot time that can be compared to good, known reference values. It is
    important to protect the integrity and authenticity of the evidence using
    some cryptographic function. In most cases, a digital signature made over
    the evidence by the \ac{TEE} itself gives proof of authenticity and is
    verified through a public key infrastructure managed by the \ac{TEE}
    manufacturer. The information included in the evidence, instead, is used to
    determine if the workload is running on an up-to-date platform and has not
    been tampered with during the provisioning phase.
\end{enumerate}

Together, these two features can give strong confidentiality, integrity and
authenticity guarantees to an application running within a \ac{TEE}. Other
mechanisms offered by some \acp{TEE} can additionally provide some support for
secure storage (e.g., \emph{data sealing}) and secure I/O. In contrast,
availability is typically out of scope of \ac{TEE} designs. Even though some
academic proposals do exist to provide some form of trusted
scheduling~\cite{alder_2021_aion}, commercial \acp{TEE} do not give any
guarantees about the execution of applications, which is still managed via
traditional means by the untrusted operating system or hypervisor.

Another important property of \acp{TEE} is a significant reduction of the
\ac{TCB}, since large software components such the host operating system and the
hypervisor are kept outside of the secure environment, thus leaving a much
smaller attack surface. This ``breaks'' the layered model of traditional systems
(cf.~\cref{thesis:intro-concepts}), allowing to isolate core workloads from the
rest of the software stack. Besides, part of the \ac{TCB}, such as application
code and \ac{TEE} version, can be verified using remote attestation to ensure
that they are in an expected state, giving additional guarantees about the
current state of the workload. Here, it is worth mentioning that the entire
attestation infrastructure is also included in the \ac{TCB}.

\subsection{Hardware-based Isolation}

\begin{figure}[t]
    \centering
    \includegraphics[width=0.8\columnwidth]{tee.drawio.pdf}
    \caption[Main TEE architectures]{Main \acs{TEE} architectures. The colored
    boxes represent trusted elements that are part of the \acs{TCB}. Software
    components, as well as some parameters related to the \ac{TEE} hardware such
    as the firmware version, can be verified via remote attestation.}
    \label{fig:phd-intro-tees}
\end{figure}

Existing \acp{TEE} can be divided in two main categories:
%
\begin{paraenum}
    \item \emph{Process-based}, where the isolation happens at the application
    layer, while the underlying operating system remains untrusted, and
    \item \emph{\acs{VM}-based}, where an entire virtual machine is protected,
    including guest firmware and operating system, while the host hypervisor is
    untrusted.
\end{paraenum}
%
Both architectures are illustrated in \cref{fig:phd-intro-tees} and discussed in
more detail below.

\paragraph{Process-based TEEs}

\iffalse
\begin{figure}[t]
    \centering
    \includegraphics[width=0.8\columnwidth]{process-based.drawio.pdf}
    \caption[Process-based TEEs]{Process-based \acsp{TEE}.}
    \label{fig:phd-intro-process-based}
\end{figure}
\fi

Process-based \acp{TEE} allow to isolate a single application, or part of it,
inside an isolated environment that is commonly known as \emph{enclave}.
Enclaves are loaded in a protected memory region that is inaccessible by neither
the ``untrusted'' application nor the operating system. Instead, they have a
well-defined entry point that can be called from the outside and used to request
some services, such as executing a function or storing some data in secure
memory. Likewise, enclave code can call untrusted code to, e.g., execute system
calls. In both cases, the processor performs a context switch to save or restore
enclave state. The lifecycle of enclaves is fully managed by the untrusted
application, which is responsible for calling the CPU instructions to create a
new empty enclave, load pages in memory, initialize the enclave, and finally
calling its functions.

In the industry, the most popular process-based \ac{TEE} is Intel
\ac{SGX}~\cite{sgx}, available on high-end x86 processors. On ARM devices,
TrustZone~\cite{alves:trustzone} allows partitioning memory into a ``normal''
and ``secure'' world. While not being a \ac{TEE} by itself, it can be paired
with a trusted OS such as OP-TEE~\cite{OPTEE} to provide enclave-like execution
environments in the secure world. Finally, although many academic solutions
exist, in this dissertation we only focus on Sancus~\cite{sancus}, a \ac{TEE}
designed for low-end microprocessors based on the OpenMSP430 architecture.

Process-based \acp{TEE} require applications to be partitioned into a secure and
non-secure component. To do so, most \acp{TEE} provide a \ac{SDK} to manage the
enclave lifecycle, call enclave functions, handle system calls and data transfer
between secure and non-secure code, and perform remote attestation. This gives
high flexibility and allows isolating sensitive code from the rest of the
(possibly big) codebase, which helps minimizing the \ac{TCB} and thus the attack
surface. At the same time, however, it requires a significant development effort
to carefully design enclave applications, locate and isolate sensitive
functions, manage and sanitize data that goes in and out of the enclave. This
can be prohibitively expensive, especially when redesigning existing (legacy)
applications that consist of thousands or millions lines of code, and it also
requires additional expertise in application developers. Several solutions exist
to mitigate these issues, such as frameworks like Open
Enclave~\cite{openenclave} that facilitate the development of enclave
applications by abstracting the \ac{TEE} layer away, or library OSes and
runtimes like Graphene-SGX~\cite{grapheneSgx} (now called Gramine) that provide
a middleware layer offering OS-like functionalities inside the enclave. In the
latter case, it is sometimes possible to port existing applications to enclaves
with little to no changes to the codebase, which significantly reduces the
development effort but at the same time it increases the \ac{TCB}.

\paragraph{VM-based TEEs}

Given the complexity and the development effort to port applications to
enclaves, coupled with a rising need to protect and isolate sensitive \ac{VM}
workloads in the cloud, recent \ac{TEE} architectures take a different approach,
offering isolation mechanisms not at the application but at the virtualization
layer. In this scenario, entire \acp{VM} are loaded into secure memory, which
includes guest firmware (not shown in \cref{fig:phd-intro-tees}), kernel, and
applications. In the industry, this functionality is offered by AMD
\sevsnp{}~\cite{AmdSevSnpWhitepaper}, Intel \ac{TDX}~\cite{intelTdxWhitepaper},
and ARM \ac{CCA}~\cite{armCCA}.

\acs{VM}-based \acp{TEE} do not require any changes to applications because
there is no \ac{TEE} boundary between the application layer and the guest
operating system. Thus, applications can execute system calls as usual, without
incurring a context switch to go in or out of the \ac{TEE}. Not only does this
greatly reduce the burden to port applications to \acp{TEE}, but it can also
improve performance in some cases. However, this comes at the cost of a
significantly bigger \ac{TCB} compared to enclaves.

Even though no changes are needed in the application layer, the guest \ac{VM}
still needs to be \ac{TEE}-aware in order to execute, e.g., to properly
configure page metadata to select which pages should be encrypted and which
should be shared with the hypervisor. Depending on the case, these changes can
be either implemented in the guest firmware or in the guest operating system. In
the open-source community, support for AMD \sevsnp{} and Intel \ac{TDX} is in
active development and some patches have already been upstreamed in the Linux
kernel, as well as in QEMU and \ac{OVMF}. Nevertheless, it is also possible to
run ``legacy'' (i.e., non-\ac{TEE}-aware) \acp{VM} by adding a privileged
\emph{paravisor} layer inside the \ac{VM} responsible for the correct
configuration of the workload. This is, for example, supported in AMD \sevsnp{}
via the Linux \ac{SVSM}~\cite{coconutSVSM}. Apart from compatibility support for
legacy workloads, a paravisor can also enable additional services such as a
\ac{vTPM}~\cite{narayanan2023svsm}.

In the cloud, \acs{VM}-based \acp{TEE} can offer strong hardware-based workload
isolation, which can be used either as a defense-in-depth approach to isolate
different tenants (together with hypervisor-based isolation), or as a mechanism
to prevent unauthorized access to a tenant's sensitive data by a malicious cloud
provider. In this latter scenario, workloads running in \acs{VM}-based \acp{TEE}
are commonly called \acp{CVM}.

\subsection{Remote Attestation}

\begin{figure}[t]
    \centering
    \includegraphics[width=.95\columnwidth]{rats.drawio.pdf}
    \caption[Simplified RATS architecture]{Simplified version of the \acs{RATS}
    architecture~\cite{ietfRats} applied to \acsp{TEE}. The colored boxes
    represent the elements whose trustworthiness is assessed by the verifier
    when checking the attestation report. Verifier, policy and reference values
    are proper \acs{RATS} terms, whereas for other relevant components a mapping
    to the official terminology is reported between brackets.}
    \label{fig:phd-intro-rats}
\end{figure}

Even though each \ac{TEE} implements remote attestation differently, the key
concepts are common to all technologies and have been standardized in the IETF
RFC 9334 under the \ac{RATS} architecture~\cite{ietfRats}.
\cref{fig:phd-intro-rats} shows a simplified version of this architecture.

Attestation is typically initiated by the \emph{verifier}, who requests a report
(or \emph{evidence} in \ac{RATS} terms) from the \ac{TEE} instance
(\emph{attester}). Such a report is retrieved from the \ac{TEE} runtime, e.g.,
the \ac{QE} in Intel \ac{SGX} and \ac{TDX}, or the AMD \ac{SP} in \sevsnp{}.
Freshness can be guaranteed by attaching some custom data to the report, such as
a nonce or a timestamp, which may be provided by the verifier itself. To verify
the authenticity of the report, the verifier has to retrieve the root of trust
(\emph{endorsement}) from the \ac{TEE} manufacturer (\emph{endorser}). This root
of trust can vary according to the \ac{TEE} and attestation scheme: In recent
\acp{TEE}, the root of trust is a \ac{PKI} root certificate that can be used to
verify the signature and certificate chain of the report. Older or academic
\acp{TEE} use different techniques, such as the \ac{EPID} scheme in Intel
\ac{SGX}~\cite{SGXexplained} or a symmetric key derivation in Sancus. After
verifying that the report is authentic, i.e., coming from a legitimate \ac{TEE},
the verifier checks its content against some predefined policy and reference
values to make sure that the \ac{TEE} is in the expected state, which involves
both the platform and the instance itself. Regarding the former, the verifier
checks if the \ac{TEE} firmware and runtime are up-to-date, and if the host has
the desired features (e.g., some BIOS settings are enabled or disabled).
Concerning the latter, instead, the verifier checks the launch policy, which
contains information such as whether debugging is enabled or not or if instance
migration is allowed, and finally compares the \emph{launch measurement} against
its reference value. This measurement consists of a hash over the initial memory
layout of the \ac{TEE} instance and may include additional information such as
page metadata. Successfully verifying the launch measurement gives the verifier
a guarantee that the \ac{TEE} instance is running the expected code. At this
point, the attestation process is complete and, if all checks are successful,
the \ac{TEE} instance can be trusted.

\iffalse
The \ac{RATS} architecture gives a more detailed and formal definition of the attestation
process, also including additional entities such as the relying party, i.e., the
entity to whom the verifier delivers the attestation results. Besides, it also
provides an in-depth description of the trust model, message formats, and
strategies to guarantee freshness of the attestation evidence. 
\fi

It is important to note that the attestation report only reflects the
\emph{initial state} of the running \ac{TEE} instance, i.e., the initial memory
layout and parameters. For \acs{VM}-based \acp{TEE}, as will be clear later on
in \cref{ch:cloud}, only the guest firmware is reflected in the initial
measurement, thus the verifier needs to employ additional techniques to
``extend'' this measurement to upper layers such as the guest OS and
applications. Yet, measuring binary files only guarantees integrity at boot
time, but cannot detect runtime attacks such as memory corruption. For detecting
this class of attacks, \ac{TEE} attestation is out of scope and other techniques
such as static/dynamic analysis and control-flow attestation should be used
instead~\cite{alder2024pandora,sha2024control}.

\subsection{Threat Model}
\label{thesis:intro-threat}

To better understand the concepts of this dissertation, it is important to give
a general overview of the \ac{TEE} threat model. A more detailed discussion of
assumptions and attacker capabilities is then provided separately for each
contribution, in the respective chapters (cf.~\cref{thesis:intro-outline}).

\paragraph{Isolation}

As mentioned above, \acp{TEE} give confidentiality, integrity and authenticity
of code and data running within them. This is achieved by means of hardware
extensions to enforce isolation, software libraries, tools and \acp{SDK} for the
development and management of \ac{TEE} applications, and an attestation
infrastructure that gives the possibility to verify security claims. These
elements are provided by the \ac{TEE} manufacturer and are thus \emph{trusted}
by assumption (although components that are open-source can be manually
reviewed). Furthermore, we also assume that all cryptography is secure and
cannot be broken, including encryption algorithms, hashes, digital signatures
and public key infrastructure. Instead, all other software running outside of
the \ac{TEE} such as host hypervisor, host operating system and other
applications, along with actors such as the infrastructure provider, local
tenants in the system and remote users are by design untrusted and possibly
malicious. 

In this model, \acp{TEE} can protect the confidentiality and integrity of
applications even when the underlying infrastructure is fully compromised, but
cannot help if the application itself is vulnerable. Attackers that are able to
exploit a bug in the code can easily break \ac{TEE} isolation and access or
modify data. Similarly, an adversary might try to exploit the interface between
the code running inside the \ac{TEE} and the untrusted world, e.g., sending
malicious inputs or modifying outputs. In both cases, it is often up to the
software developer to prevent such attacks from happening, e.g., by keeping the
\ac{TCB} small and sanitizing interfaces.

\paragraph{Attacks on TEEs}

Despite the isolation guarantees provided by \acp{TEE}, there are some attacks
that current technologies may be vulnerable to. First, even though \acp{TEE}
protect memory from unauthorized access, a privileged adversary can observe
access patterns and infer secret data via side
channels~\cite{lou2021survey,canella2019systematic,shen2021sidechannelsok}.
Since \acp{TEE} assume a fully compromised host, attackers can leverage
techniques to precisely control the execution of \ac{TEE} applications even at
instruction-level granularity (i.e.,
\emph{single-stepping})~\cite{controlled-channel-attacks,jo2017sgxstep,wilke2024sevstep}.
Side-channel exploitation on \acp{TEE} has been repeatedly demonstrated in the
past, especially on Intel \ac{SGX}~\cite{schaik2024sgxsok}. Recently, techniques
such as AEX-Notify have been implemented to mitigate single-stepping on
enclaves~\cite{constable2023aexnotify}; Such mechanisms, while not addressing
the underlying root cause of side channels, can make it much harder to perform
this class of attacks. However, motivated attackers might still be able to
succeed, as demonstrated by a recent attack on Intel
\ac{TDX}~\cite{tdxdownCCS24}. In general, side-channel attacks are considered
out of scope by current commercial \acp{TEE} and should be mitigated by means of
software mechanisms such as constant-time
programming~\cite{cauligi2020constanttime} and/or oblivious
computation~\cite{alam2023oblivioussok,sasy2017zerotrace}.

Second, attackers might be able to perform physical attacks to induce faults
during some sensitive computation and break the confidentiality and/or integrity
of \ac{TEE}
applications~\cite{murdock2020plundervolt,chen2021voltpillager,lipp2021platypus}.
As with side channels, software-based mitigations might also be used here to
deal with these attacks~\cite{wichelmann2024obelix}, combined with
hardware-based mechanisms to restrict the usage of certain features such as the
voltage scaling interface.

Third, the \ac{TEE} implementation itself may present design flaws and bugs that
can break the confidentiality and integrity of applications. This has been shown
by researchers on several
occasions~\cite{morbitzer2018severed,securingUnseen,cui2021smashex,schluter2024heckler,schluter2024wesee,li2022cipherleaks,schaik2024sgxsok}.
It is worth mentioning that most of these vulnerabilities do not apply to all
\acp{TEE} in general but are specific to a particular architecture. While most
of them have been already patched and can be mitigated by running applications
on up-to-date platforms (which can be verified using remote attestation), it is
also possible that new vulnerabilities are discovered in the future. If
attackers are able to exploit a zero-day vulnerability on a particular
architecture, they might be able to break the isolation provided by the
\ac{TEE}. Unfortunately, there is little that administrators and application
developers can do to deal with these attacks, apart from applying traditional
defense-in-depth techniques and reducing the attack surface as much as possible.

All the aspects mentioned above are crucial when deploying a \ac{TEE} workload
in the real world, but they are out of the scope of this dissertation.
Nevertheless, we briefly touch upon them in our concluding chapter
(cf.~\cref{thesis:concl-future}).